% Generated draft: requires mathematical and source review.
\documentclass[11pt]{article}
\usepackage[margin=25mm]{geometry}
\usepackage{fontspec}
\setmainfont{Latin Modern Roman}
\usepackage{amsmath,amssymb,mathtools,unicode-math}
\setmathfont{Latin Modern Math}
\usepackage{enumitem,needspace,xurl,hyperref}
\setlength{\parindent}{0pt}
\setlength{\parskip}{5pt}
\setlength{\emergencystretch}{3em}
\newcommand{\N}{\mathbb{N}}
\newcommand{\Z}{\mathbb{Z}}
\newcommand{\Q}{\mathbb{Q}}
\newcommand{\R}{\mathbb{R}}
\newcommand{\C}{\mathbb{C}}
\newcommand{\F}{\mathbb{F}}
\newcommand{\A}{\mathbb{A}}
\newcommand{\PP}{\mathbb{P}}
\newcommand{\E}{\mathbb{E}}
\newcommand{\eps}{\varepsilon}
\newcommand{\dd}{\,\mathrm{d}}
\begin{document}
\section*{UnsolvedMath ID: 30003075}
\textbf{Source problem code:} catalogue entry
\subsection{Definitions and mathematical statement}
The supplied question asks for a characterization of finite point configurations producing unexpected curves. It does not specify the ambient space, ground field, characteristic, curve degree, or fat-point multiplicity. The following gives a precise conditional plane-curve formulation; these additional parameters are not asserted to be hypotheses supplied by the source.

Suppose the intended ambient space is \(\PP^2_k\), where \(k\) is a specified algebraically closed field. Let \(Z\subset\PP^2_k\) be a finite set of distinct points, regarded as a reduced subscheme, and let \(\mathcal I_Z\) be its ideal sheaf. For integers \(d\geq 0\) and \(m\geq 1\), and a point \(P\notin Z\), define the fat point \(mP\) by the ideal sheaf \(\mathcal I_P^m\). Set
\[
V_d(Z)=H^0\bigl(\PP^2_k,\mathcal I_Z(d)\bigr),\qquad
V_{d,m}(Z,P)=H^0\bigl(\PP^2_k,(\mathcal I_Z\cap\mathcal I_P^m)(d)\bigr).
\]
Thus \(V_{d,m}(Z,P)\) consists of degree-\(d\) homogeneous forms vanishing on \(Z\) and to order at least \(m\) at \(P\). The restriction map
\[
\rho_{Z,P}^{d,m}:V_d(Z)\longrightarrow H^0\bigl(mP,\mathcal O_{mP}(d)\bigr)
\]
has kernel \(V_{d,m}(Z,P)\), and its target has dimension \(\binom{m+1}{2}\). The expected vector-space dimension is therefore
\[
e_{d,m}(Z)=\max\left\{0,\dim_k V_d(Z)-\binom{m+1}{2}\right\}.
\]
For these specified parameters, define the unexpectedness condition by
\[
\dim_k V_{d,m}(Z,P)>e_{d,m}(Z)
\quad\text{for general }P\in\PP^2_k\setminus Z.
\]
Here ``general'' means that the inequality holds for every \(P\) in some nonempty Zariski-open subset of \(\PP^2_k\setminus Z\). Equivalently, the restriction map fails to have maximal rank for general \(P\).

\textbf{Requested characterization, subject to unresolved scope.} Characterize the finite point sets \(Z\) satisfying this unexpectedness condition in the ambient setting and degree--multiplicity regime intended by the source. The supplied record does not determine whether \(d,m\) are fixed, whether existence of a pair \((d,m)\) is intended, or whether a prescribed relation between them is required. It also does not establish that the conditional plane-curve setting above is the complete intended setting. Consequently, a unique fully specified classification problem cannot be recovered from the supplied wording alone.
\subsection{Short English statement}
Characterize finite point configurations for which requiring a curve through the configuration to have a specified multiplicity at a general additional point leaves more curves than the expected dimension predicts. A precise plane-curve interpretation is given conditionally; the source leaves the ground field, ambient space, and degree--multiplicity quantifiers unspecified.
\subsection{Sources}
\begin{itemize}
\item Local UnsolvedMath registry, supplied complete record: permanent ID \(30003075\), title \emph{Point Configurations Producing Unexpected Curves}, exact source problem code \texttt{catalogue entry}, in the Oberwolfach Reports Open Problems collection. The record gives the proposed year as 2016 and supplies no individual proposer.
\item Supplied primary provenance: \emph{Mini-Workshop: Arrangements of Subvarieties, and their Applications in Algebraic Geometry} (2016), Oberwolfach Reports, pp.~653--702, DOI \texttt{10.4171/owr/2016/14}. Supplied source URL and literature-source label ``OWR report 14/2016'': \url{https://doi.org/10.4171/owr/2016/14}.
\item Supplied literature source: Cook--Harbourne--Migliore--Nagel, \url{https://doi.org/10.1112/S0010437X18007376}. The supplied record does not give a title or a specific theorem reference.
\item Supplied literature source: Dimca--Sticlaru, \url{https://doi.org/10.1007/s10711-021-00602-5}. The supplied record does not give a title or a specific theorem reference.
\item Literature assessment date: \texttt{2026-08-21}. This is supplied dated curation, not independently verified. The supplied status is ``partially solved''. Its assessment says that the broad classification question remains open, that an unexpected-curve criterion is characterized through the splitting type of the dual line arrangement, and that many explicit families and low-degree cases are known. The background describes this as dated literature triage derived from the supplied Oberwolfach curation and calls for independent verification of time-sensitive status and cited claims.
\item The supplied statement-status label is \texttt{corrected\_verified}; its accompanying qualification says only that notation, typography, and sentence structure were repaired without changing mathematical content. This label is not treated here as independent mathematical verification.
\end{itemize}
\end{document}
